% !TeX root = 控制工程基础笔记.tex
% ============================================================
%  封面：李思风（极简学术）—— 白底、居中标题、无装饰边框
%  参照《线性代数：原理与应用》（李思，清华大学）封面版式
%
%  可用字段（在主文件导言区 \newcommand 预定义即可覆盖）：
%    \notename      主标题（必填，主文件定义）
%    \noteauthor    作者（必填，主文件定义）
%    \notesubtitle  标题下方小字，默认 \today（日期/版本行）
%    \noteextra     底部单位/学院行，默认不显示
%
%  说明：中部插图为「子昕手绘」（由 figures/prep_cover_image.m 预处理，
%        见 figures/子昕手绘_封面.png），整段 \includegraphics 删除后
%        即为纯文字极简封面。
% ============================================================
\begin{titlepage}
	\thispagestyle{empty}
	\providecommand{\notesubtitle}{\today}
	\providecommand{\noteextra}{}
	\begin{center}
		\vspace*{3.8cm}

		% ---- 主标题 ----
		{\bfseries\songti\fontsize{24.8}{30}\selectfont \notename\par}

		\vspace{3.0cm}

		% ---- 日期 / 版本行 ----
		{\songti\fontsize{14.4}{18}\selectfont （\notesubtitle）\par}

		\vspace{2.4cm}

		% ---- 封面插图：子昕手绘（可整段删除）----
		\includegraphics[width=5.0cm]{figures/子昕手绘_封面.png}\\[0.6em]
		{\fangsong\fontsize{10.5}{12}}

		\vspace{2.6cm}

		% ---- 作者（楷体）----
		{\kaishu\fontsize{24.8}{30}\selectfont \noteauthor\par}

		\vspace{1.5cm}

		% ---- 单位 / 学院（未定义则不显示）----
		{\songti\fontsize{20.7}{26}\selectfont \noteextra\par}

		\vspace*{1.4cm}
	\end{center}
\end{titlepage}
